

\section*{Preface} 

This book started out as a collection of lecture notes for a course on transducers, taught at the University of Warsaw in 2024 and 2026, but it is meant to evolve into a more comprehensive resource on transducers.  I would recommend visiting the online version of the book, which links to a Lean formalisation of the entire text.  I hope that the presence of such a formalisation will allow others to build on results presented here without fear of errors. Also, I hope that the Lean formalisation can free me to adopt a readable and informal writing style, with pictures and examples, since doubts can be resolved conclusively by reading the linked formalisation.

\paragraph*{Acknowledgements.} I would like to thank the students of the course for their feedback, and Antek Puch for his support in selecting exercises. I would also like to thank Rafał Stefański and Anshul Goyal for helping start the Lean formalisation project. Ultimately, the formalisation was done using the AI assistant Aristotle. (I confess that I have done  little more than spot checks to see if the formalisation crresponds to the text, but I hope that convenient web interface with direct links to Lean will ensure that any confabulations are exposed.)  I have also used AI, mainly Claude, to find errors in the text, and to create the website. The website is based on reflowtex by Radek Piórkowski, which allows displaying TeX on web pages.

\paragraph*{Future work.} Here are some plans for future editions of this book.
\begin{itemize}
    \item Now Lean covers only parts A-C, but D should be done soon.
    \item Not all exercises have solutions, and these solutions are not formalised in Lean. Eventually, all exercises should have solutions with Lean formalisations.
    \item Regular list functions and \mso interpretations.
\end{itemize}
Further down the line, I would like to extend from strings to trees and graphs.